\SourceLeaf{376}%
purposes. By about 1811 a triangulation network from Fiume to Turin had been completed. It thus extended over seven degrees of longitude along a parallel at the 45th degree of latitude. To put this work to broader scientific use in determining the shape of the Earth, the French government of the time decided, at \textsc{Laplace}’s suggestion, to extend the triangulation network westward to the Atlantic Ocean and to combine it with the operations needed to measure an arc of longitude. The work began at once and with enthusiasm. Events interrupted it for a time, but it soon resumed with equal vigour. By 1818 the triangulation network had been surveyed across French territory from the Atlantic Ocean near Bordeaux to the border of Savoy. The only part of the geodetic work still to be completed was the section in the dominions of the king of Sardinia. The Sardinian and Austrian governments, both keenly aware of the scientific importance of this great undertaking, decided to appoint a commission of astronomers and officers from both states to carry out the remaining geodetic operations and the astronomical operations needed in Italy. These operations are the subject of the present work, which appears to have been jointly edited by the astronomers \textsc{Carlini} and \textsc{Plana}.

The first volume is devoted entirely to the geodetic operations. The two eastern endpoints of the triangulation network in France, Mont Colombier and Mont Granier (near Chambery), form the side from which the new survey had to start and continue to the westernmost side of the network in Lombardy, Massé --- Superga (near Turin). One might have expected that in this terrain, where the highest mountains in Europe are to be found, it would be easy to form large triangles and that a very small number of intermediate points would suffice to make the connection --- Mont Granier is only 150000 metres from Superga. Yet the very opposite was true. This modest stretch presented the greatest difficulties, since the summits of the higher mountains are inaccessible or difficult to reach, the materials for building the signals can be carried up only with the greatest effort, and violent storms both
